\section{Mechanism evidence, limits and conclusion}
\label{sec:7}
On observed Confirm96-B, an offline oracle finds a modest average native
advantage in local ridge-gradient compatibility at a common query-only state.
Its finite-amplitude loss response is linked to that geometry by construction.
Neither diagnostic establishes a positive target-level association with trained
folding gains. An additional Train8-derived finite parameter displacement does
not establish positive mean transfer or native superiority on Dev8. The earlier
shared differential-response numerical gate remains failed. Full methods,
negative results and precision checks are retained in the appendix; these
experiments do not provide a deployable inference procedure.

A descriptive decomposition locates 75.2\% of the Train384 Protenix improvement
in reference pairs separated by at least 24 sequence positions. These pairs
constitute 57.40\% of eligible pairs and also show the largest within-group gain
(+0.05148). Different training settings have different score-headroom patterns.
This shows where distance preservation improves, not recovery of coevolution
or a universal long-range specialization.

The central evidence concerns particular checkpoints, interfaces, training
recipes and target distributions. Mean-preserving rotations do not exhaust
all direction controls. Repeatedly used targets and few training seeds limit
inference beyond the fitted systems. Sequence filtering does not guarantee
family or pretraining isolation, and the incomplete historical CATH audit does
not cover all Train384 exposure. The AF2 results explicitly separate direction
sensitivity from the competitive value of the factor constraint. Pending
AtlasFold adaptation and system-only baselines do not strengthen that claim
until their corresponding comparisons are complete.

\textbf{Conclusion.} Native operator-centered residuals provide a concrete way to
adapt frozen structure predictors. Trained orthogonal controls show that
orientation relative to frozen downstream computation can affect structural
quality beyond the factor form and decoder spectrum. This effect extends
across tested settings and to AF2, but depends on training and target conditions;
it does not ensure superiority over generic adaptation. Fixed Protenix adapters
also retain useful updates beyond their adaptation-training length range.
These findings define both the value and the limits of native update alignment.
